% Cascada genética: de las condiciones a lo condicionado. `genesis` empuja hacia abajo;
% la fuerza se rotula con el nombre del autor. Niveles al margen.
\documentclass[tikz,border=12pt]{standalone}
\usepackage{diagrama}
\begin{document}
\begin{tikzpicture}[every node/.append style={font=\px{30}}]
  \node[titulo=44] at (3,7) {Título de la génesis};
  \node (n1) at (3,5.6) {Condición};
  \node[caja, minimum width=2.6cm] (n2) at (1,3.6) {Facultad 1};
  \node[caja, minimum width=2.6cm] (n3) at (5,3.6) {Facultad 2};
  \node (n4) at (1,1.6) {Objeto 1};
  \node (n5) at (5,1.6) {Objeto 2};
  \draw[genesis] (n1) to[out=-150, in=90] (n2);                               % e1
  \draw[genesis] (n1) to[out=-30, in=90] (n3);                                % e2
  \draw[genesis] (n2) -- node[rotulo] {\textit{Fuerza}} (n4);                 % e3
  \draw[genesis] (n3) -- node[rotulo] {\textit{Fuerza}} (n5);                 % e4
  \draw[reciproca rayada] (n4) -- node[rotulo] {\textit{Correlación}} (n5);   % e5
  \node[nivel=20, anchor=west] at (6.8,5.6) {VIRTUAL};
  \node[nivel=20, anchor=west] at (6.8,3.6) {FACULTADES};
  \node[nivel=20, anchor=west] at (6.8,1.6) {ACTUAL};
\end{tikzpicture}
\end{document}
